% from intro-fr.tex
\oldpage[I]{6}
n'essayerons pas non plus de donner ici une introduction <<~intuitive~>> aux notions développées dans le chapitre premier.
Le lecteur qui désirerait avoir un aperçu préliminaire des matières de ce Traité pourra se reporter à la conférence faite par A.~Grothendieck au Congrès international des Mathématiciens à Edinburgh en 1958~\cite{I-7}, et à l'exposé~\cite{I-8} du même auteur.
Le travail~\cite{I-14} (cité~(FAC)) de J.-P.~Serre peut aussi être considéré comme un exposé intermédiaire entre le point de vue classique et le point de vue des schémas en Géométrie algébrique, et à ce titre, sa lecture peut constituer une excellente préparation à celle de nos \emph{Éléments}.

\bigskip
\begin{center}*\quad*\quad*\end{center}
\bigskip

A titre informatif, nous donnons ci-dessous le plan général prévu pour ce Traité, d'ailleurs sujet à modifications ultérieures, surtout en ce qui concerne les derniers chapitres~:

\medskip
\begin{tabular}{rl}
  Chapitre Premier. & --- Le langage des schémas.\\
  --- \textbf{II}. & --- Étude globale élémentaire de quelques classes de morphismes.\\
  --- \textbf{III}. & --- Cohomologie des faisceaux algébriques cohérents. Applications.\\
  --- \textbf{IV}. & --- Étude locale des morphismes.\\
  --- \textbf{V}. & --- Procédés élémentaires de construction de schémas.\\
  --- \textbf{VI}. & --- Technique de descente. Méthode générale de construction des schémas.\\
  --- \textbf{VII}. & --- Schémas de groupes, espaces fibrés principaux.\\
  --- \textbf{VIII}. & --- Étude différentielle des espaces fibrés.\\
  --- \textbf{IX}. & --- Le groupe fondamental.\\
  --- \textbf{X}. & --- Résidus et dualité.\\
  --- \textbf{XI}. & --- Théories d'intersection, classes de Chern, théorème de Riemann-Roch.\\
  --- \textbf{XII}. & --- Schémas abéliens et schémas de Picard.\\
  --- \textbf{XIII}. & --- Cohomologie de Weil.\\
\end{tabular}
\medskip

En principe, tous les chapitres sont considérés comme ouverts, et des paragraphes supplémentaires pourront toujours leur être ajoutés ultérieurement; de tels paragraphes paraîtront en fascicules séparés, pour diminuer les inconvénients du mode de publication adopté.
Lorsqu'un tel paragraphe est prévu ou en préparation au moment de la publication d'un chapitre, il sera mentionné dans le sommaire dudit chapitre, même si en raison de certains ordres d'urgence sa publication effective devait être nettement postérieure.
Pour la commodité du lecteur, nous donnons dans un <<~Chapitre~0~>> des compléments divers d'Algèbre commutative, d'Algèbre homologique, de Théorie des faisceaux, utilisés au cours des chapitres de ce Traité, qui sont plus ou moins bien connus, mais pour lesquels il n'a pas été possible de donner des références commodes.
Il est recommandé au lecteur de ne se reporter au chapitre~0 qu'en cours de lecture du Traité proprement dit, et dans la mesure où les résultats auxquels nous référons ne lui sont pas

